% Edited conversion of source pp. 43--45.
% Source plots are represented by compact TikZ/PGFPlots reconstructions.
\ReportHeading
  {Набігання ступінчастої ударної хвилі під кутом $45^\circ$ на модель типу циліндр--півсфера}
  {А.~С. Колядюк}
  {Інститут механіки ім. С.~П. Тимошенка НАН України, Київ, Україна}
  {}

Для аналізу динаміки складних оболонкових конструкцій необхідно визначати навантаження, яке безпосередньо діє на їхню поверхню. Розглянуто набігання ступінчастої ударної хвилі з крутим фронтом і стрибком тиску $\Delta P=10^6$~Па під кутом $45^\circ$ на модель типу циліндр--півсфера, що відтворює геометрію обтічника. Радіуси циліндричної та сферичної частин становлять 324~мм.

\begin{center}
\begin{tikzpicture}[x=0.58cm,y=0.58cm,line cap=round,line join=round]
  \draw[thick] (0,0)--(6.0,0)--(6.0,2.2) arc (-90:90:2.2)--(0,4.4)--cycle;
  \draw[step=0.35,very thin] (-0.2,-0.2) grid (8.5,4.6);
  \draw[white,line width=1.1pt] (0,0)--(6.0,0)--(6.0,2.2) arc (-90:90:2.2)--(0,4.4)--cycle;
  \draw[thick] (0,0)--(6.0,0)--(6.0,2.2) arc (-90:90:2.2)--(0,4.4)--cycle;
  \foreach \x in {0.35,0.7,...,5.95} {\draw[very thin] (\x,0)--(\x,4.4);}
  \foreach \y in {0.35,0.7,...,4.2} {\draw[very thin] (0,\y)--(6.0,\y);}
  \draw[->,very thick] (-1.0,5.4)--(1.2,3.2) node[midway,above right] {ударна хвиля};
\end{tikzpicture}
\captionof{figure}{Схематична скінченно-елементна сітка розрахункової області біля комбінованої циліндрично-сферичної оболонки.}
\end{center}

Середовище моделювали як нев’язкий стисливий ідеальний газ. У цій постановці усереднені за Рейнольдсом рівняння Нав’є--Стокса зводяться до рівнянь Ейлера:
\begin{align}
 \frac{\partial \rho}{\partial t}+\nabla\!\cdot(\rho\mathbf v)&=0,\\
 \frac{\partial(\rho\mathbf v)}{\partial t}
 +\nabla\!\cdot\bigl(\rho\mathbf v\otimes\mathbf v+p\mathbf I\bigr)&=0,\\
 \frac{\partial(\rho E)}{\partial t}
 +\nabla\!\cdot\bigl[\mathbf v(\rho E+p)\bigr]&=0,
\end{align}
де $\mathbf v$ --- вектор швидкості, $\rho$ --- густина, $p$ --- статичний тиск, а повна енергія на одиницю маси
\[
 E=h-\frac{p}{\rho}+\frac{|\mathbf v|^2}{2},
\]
$h$ --- ентальпія.

Методом скінченних елементів визначено дифракційний тиск на зовнішній поверхні комбінованої оболонки. Поширення хвилі моделювали на інтервалі $9\times10^{-4}$~с. Розрахунки показують різке зростання тиску під час проходження фронту, подальше затухання та появу відкладених локальних максимумів у різних точках периметра.

\begin{center}
\begin{tikzpicture}
\begin{axis}[
 width=0.84\textwidth,height=4.6cm,
 xlabel={$t$, $10^{-4}$ с},ylabel={$p_d$, МПа},
 xmin=0,xmax=9,ymin=0,ymax=6,
 xtick={0,1,...,9},grid=major,
 legend style={font=\scriptsize,at={(0.98,0.96)},anchor=north east}]
 \addplot[thick,domain=0:9,samples=180] {x<0.6 ? 0 : 5.6*exp(-(x-0.6)/1.3)};
 \addplot[dashed,thick,domain=0:9,samples=180] {x<0.8 ? 0 : 5.1*exp(-(x-0.8)/1.5)};
 \addplot[dashdotted,thick,domain=0:9,samples=180] {x<1.9 ? 0 : 3.6*exp(-(x-1.9)/1.7)};
 \addplot[dotted,thick,domain=0:9,samples=180] {x<2.8 ? 0 : 1.5*exp(-(x-2.8)/2.0)};
 \legend{передня точка,верхня дуга,бічна ділянка,тильна ділянка}
\end{axis}
\end{tikzpicture}
\captionof{figure}{Схематичне відтворення характерних залежностей дифракційного тиску від часу у вибраних точках периметра.}
\end{center}

Максимальний тиск виникає на навітряній частині оболонки й зменшується до тильної ділянки. Час досягнення максимуму зростає, відображаючи послідовне проходження фронту та дифракцію хвилі.

\begin{center}
\begin{tikzpicture}[x=0.62cm,y=0.62cm,line cap=round,line join=round]
  \draw[thick] (-3,0)--(-3,2.2) arc (180:0:3 and 3)--(3,0);
  \draw[thick] (-1.65,0)--(-1.65,1.2) arc (180:0:1.65 and 1.65)--(1.65,0);
  \foreach \a/\r/\txt in {180/3.8/5.52,150/3.5/4.81,120/3.3/4.10,90/3.1/4.92,60/2.8/3.72,35/2.5/1.97,20/2.2/1.06,8/2.0/0.55} {
    \draw ({1.65*cos(\a)},{1.65*sin(\a)})--({\r*cos(\a)},{\r*sin(\a)});
    \node[font=\scriptsize] at ({(\r+0.35)*cos(\a)},{(\r+0.35)*sin(\a)}) {$\txt$};
  }
  \node[font=\scriptsize] at (0,-0.45) {$p_{d,\max}$, МПа};
\end{tikzpicture}
\captionof{figure}{Схематична епюра максимального дифракційного тиску вздовж периметра оболонки.}
\end{center}

Отримані параметри навантаження можуть бути використані у подальшому аналізі динаміки багатошарових оболонкових конструкцій.

{\small
\begin{thebibliography}{9}
\bibitem{shock:handbook2001}
G.~Ben-Dor, O.~Igra, and T.~Elperin,
\emph{Handbook of Shock Waves}, vol.~1.
Amsterdam: Academic Press--Elsevier, 2001, 889~p.

\bibitem{shock:zeldovich2012}
Ya.~B. Zel’dovich and Yu.~P. Raizer,
\emph{Physics of Shock Waves and High-Temperature Hydrodynamic Phenomena}.
Courier Corporation, 2012, 944~p.
\end{thebibliography}
}
